\chapter{Max Flow}

\begin{definition}[Flow Network]
	A \textbf{flow network} is a directed graph $G = (V, E, c)$ with two special vertices $s$ and $t$ called the \textbf{source} and \textbf{sink}, respectively. $c: E \to \mathbb{Z}^+$ is a function that assigns a non-negative integer capacity to each edge.
\end{definition}

\begin{definition}[Flow]
	A \textbf{flow} in a flow network $G$ is a function $f: V\times V \to \mathbb{Z}$ that satisfies the following conditions:
	\begin{itemize}
		\item \textbf{Capacity constraint}: For every pair of vertices $(u, v) \in V\times V$, we have $f(u, v) \leq c(u, v)$ (let $c(u, v) = 0$ if $(u, v) \notin E$).
		\item \textbf{Skew symmetry}: For every pair of vertices $(u, v) \in V\times V$, we have $f(u, v) = -f(v, u)$.
		\item \textbf{Flow conservation}: For every vertex $u \in V\setminus\{s, t\}$, we have $\sum_{(u,x) \in E} f(u, x) = 0$. Here we implicitly assume for every edge, there's a reverse edge with capacity 0 (if it doesn't exist in the original graph).
	\end{itemize}
	$v(f) = \sum_{(s, x) \in E} f(s, x) = \sum_{(x, t) \in E} f(x, t)$ is called the \textbf{value} of the flow $f$.
\end{definition}

\section{Ford-Fulkerson Algorithm}

We give a naive algorithm to find a maximum flow:

\begin{enumerate}
	\item Find a path from $s$ to $t$.
	\item Push as much flow as possible along this path.
	\item Find another path that has \textbf{free capacity}.
	\item Push as much flow as possible along this path.
	\item Repeat steps 3 and 4 until there is no path from $s$ to $t$ with free capacity.
\end{enumerate}

The \textbf{free capacity} here seems to be a bit vague, so we need to define it more formally, in a way that makes this algorithm work:

\begin{definition}[Residual Network]
	Given a flow network $G = (V, E, c)$ and a flow $f$, the \textbf{residual network} $G_f$ is defined as follows:
	\begin{itemize}
		\item $c_f(u,v)$ represents the amount of additional flow that could be sent through $(u,v)$ without violating the capacity constraint.
		\item $c_f(u, v) = c(u, v) - f(u, v)$ for each edge $(u, v) \in E$.
		\item Only edges with positive capacity appear in $G_f$.
	\end{itemize}
\end{definition}

Now we can restate the algorithm as follows:

\begin{algorithm}[H]
	\caption{Ford-Fulkerson Algorithm}
	$G_f \gets G$\;
	Initialize $f$ to be the zero flow\;
	\While{there is a path $P$ from $s$ to $t$ in $G_f$}
	{
		$f' \coloneqq $ maximum flow along $P$\;
		$f \gets f + f'$\;
	}
	return $f$\;
\end{algorithm}

\begin{lemma}
	$f'' = f + f'$ is a legal flow in $G$.
\end{lemma}

\begin{proof}
	We need to verify the three conditions in the definition of flow:
	\begin{itemize}
		\item \textbf{Capacity constraint}: $f''(u,v) = f(u,v)+ f'(u,v)\le f(u,v) + c_f(u,v) = f(u,v) + (c(u,v) - f(u,v)) = c(u,v)$ for every pair of vertices $(u, v) \in V\times V$.
		\item \textbf{Skew symmetry}: $f''(u, v) = -f''(v, u)$ because we maintain this property throughout the algorithm. When we push flow along a path, we update the flow in both directions accordingly.
		\item \textbf{Flow conservation}: For every vertex $u \in V\setminus\{s, t\}$, we have $\sum_{(u,x) \in E} f''(u, x) = \sum_{(u,x) \in E} (f(u, x) + f'(u, x)) = \sum_{(u,x) \in E} f(u, x) + \sum_{(u,x) \in E} f'(u, x) = 0 + 0 = 0$.
	\end{itemize}
\end{proof}

\section{Max Flow-Min Cut Theorem}

To proof the correctness of the algorithm, we need to introduce the following definition:

\begin{definition}[$s$-$t$ Cut]
	$(A,B)$ is an \textbf{$s$-$t$ cut} if $A\cup B = V$, $A\cap B = \emptyset$, $s \in A$ and $t \in B$.

	The capacity of the cut is defined as $c(A, B) = \sum_{(u, v) \in E, u \in A, v \in B} c(u, v)$. For any flow $f$, $v(f) \leq c(A, B)$.
\end{definition}

\begin{theorem}[Max Flow-Min Cut Theorem]
	The following statements are equivalent:
	\begin{itemize}
		\item $f$ is a maximum flow in $G$.
		\item There is no path from $s$ to $t$ in the residual network $G_f$.
		\item There exists an $s$-$t$ cut $(A, B)$ such that $v(f) = c(A, B)$.
	\end{itemize}
\end{theorem}
\begin{proof}
	(1) $\Rightarrow$ (2): We prove the contrapositive. If there is a path from $s$ to $t$ in the residual network $G_f$, we can push more flow along this path, which contradicts the assumption that $f$ is a maximum flow.

	(2) $\Rightarrow$ (3): Let $A = \{v \in V : v \text{ is reachable from } s \text{ in } G_f\}$ and $B = V\setminus A$. Then $(A, B)$ is an $s$-$t$ cut. Since there is no path from $s$ to $t$ in $G_f$, we have $t \in B$. For every edge $(u, v)$ with $u \in A$ and $v \in B$, we have $c_f(u, v) = 0$, which implies that $f(u, v) = c(u, v)$. Therefore, $v(f) = \sum_{(s, x) \in E} f(s, x) = \sum_{u\in A} \sum_{(u, x) \in E} f(u, x) =\sum_{(u, v) \in E, u \in A, v \in B} f(u, v) = \sum_{(u, v) \in E, u \in A, v \in B} c(u, v) = c(A, B)$.

	(3) $\Rightarrow$ (1): For any flow $f'$, we have $v(f') \leq c(A, B) = v(f)$, which implies that $f$ is a maximum flow.
\end{proof}

The efficiency of the Ford-Fulkerson algorithm is $O(m \cdot v(f))$, as it takes $O(m)$ time to find a path, and augmenting the flow along the path can increase the flow value by at least 1. However, if the capacities are irrational numbers, the algorithm may not terminate.

To improve Ford-Fulkerson Algorithm, we can choose path greedily, so that the bottleneck edges are maximized.

\begin{lemma}
	Denote the next augmenting path as $p$, and let $\delta := v(f_{\max})-v(f)$ be the value of max flow in $G_f$. Then we have $v(f_p) \ge \delta/m$ if we maximize bottleneck edges.
\end{lemma}
\begin{proof}
	Here we proof by contradiction. Assume that $v(f_p) < \delta/m$.

	Then we can remove all edges with $c_f(e) < \delta/m$ from $G_f$, and expect $s$ and $t$ are unconnected.

	The capacity of such cut is less than $m \cdot \delta/m = \delta \Rightarrow \contra$.
\end{proof}

Therefore, $v(f_p)\ge \delta/m = (v(f_{\max})-v(f))/m$.
If we denote $\Delta_i := v(f_{\max})-v(f_i)$, where $f_i$ is the flow after the $i$-th iteration, we have $\Delta_{i+1} \leq \Delta_i - \frac{\Delta_i}{m} = \Delta_i (1 - \frac{1}{m})$. This implies that $\Delta_i \leq v(f_{\max})(1 - \frac{1}{m})^i < v(f_{\max}) e^{-i/m}$. To guarantee that $\Delta_i < 1$, $i = m \log v(f_{\max})$ suffices for integer-capacity graphs.

It means that the number of iterations is at most $O(m \log v(f_{\max}))$, so the total time complexity is $O(m^2 \log v(f_{\max}))$.

\section{Examples using Max Flow}

\begin{question}[Max Matching in Bipartite Graph]
	Given a bipartite graph $G = (V, E)$ with $V = L \cup R$, find the maximum matching in $G$.
\end{question}

One can use it to prove Hall's theorem:
\begin{theorem}[Hall's Theorem]
	If a bipartite graph $G = (L \cup R, E)$ does not have a perfect matching, then there exists a subset $S \subseteq V$ such that $|N(S)| < |S|$, where $N(S)$ is the set of neighbors of $S$ in $V$.
\end{theorem}
